
où les
flèches horizontales sont définies en envoyant @@M0@@ sur @@M1@@ pour @@M2@@, et @@M3@@ sur @@M4@@ pour @@M5@@, où l'homomorphisme @@M6@@ est défini en envoyant @@M7@@ sur @@M8@@ pour @@M9@@, et @@M10@@ sur @@M11@@ pour @@M12@@, et l'homomorphisme @@M13@@ en annulant les @@M14@@, @@M15@@. On a alors:
@@M16@@
La situation étudiée se déduit donc de la situation @@M17@@ par le changement de base @@M18@@. Or les immersions correspondantes @@M19@@, @@M20@@ et @@M21@@ sont régulières; par suite, les groupes d'homologie des complexes de Koszul associés aux suites @@M22@@, @@M23@@ et @@M24@@, sont respectivement les @@M25@@, @@M26@@, @@M27@@. Par suite:
@@M28@@

On utilise alors la suite spectrale de transitivité ([1], XVI 5)
@@M29@@

D'après les hypothèses de i), @@M30@@ pour @@M31@@ ou @@M32@@. Donc elle dégénère, et @@M33@@ est régulière d'après (1.7.2).

D'après (1.6.1), la suite:
@@M34@@
est exacte. Comme @@M35@@ est régulière, les trois faisceaux conormaux sont localement libres de type fini; de plus, leur rang en un point est le
